\documentclass[12pt,a4paper]{article}

\usepackage[utf8]{inputenc}
\usepackage[T1]{fontenc}
\usepackage[portuguese]{babel}

\usepackage{listings}
\usepackage{xcolor}
\usepackage{amsmath}
\usepackage{amssymb}
\usepackage{amsfonts}

\lstset{
    basicstyle=\ttfamily\small,
    keywordstyle=\color{blue}\bfseries, 
    commentstyle=\color{green!70!black},
    stringstyle=\color{red}, 
    numbers=left, 
    numberstyle=\tiny\color{gray}, 
    frame=single, 
    breaklines=true, 
    showstringspaces=false 
}

\usepackage[margin=2.5cm]{geometry}

\begin{document}

\begin{center}
    {\LARGE MC521 - Relatório I} \\ [2,5em]
    {\large Júlio da Silva Telles} \\ [1,0em]
    {\large 12 de Julho de 2026} \\ [1,5em]
\end{center}

\vspace{1cm} 

\section{F - P06 (AtCoder - abc251\_d)}

O problema F descreve uma situação em que dado um inteiro $w$ tal que $1 \le w \le 10^6$ e necessario encontrar uma
familia de elementos descrita como $W = \{ w_i \in \mathbb{Z} \}_{i \in I} \text{ com } 1 \le|I|
\le 300, \text{ tal que } \forall n \in \mathbb{Z} \cap (1, w), \, \exists A \subseteq I 
\text{ t.q. } 1 \le |A| \le 3 \text{ e } n = \sum_{a \in A} w_a$. 

\subsection{Solução}

Como w não pode ser maior que $10^6$ e não existe limitações adicionais ao problema podemos apenas formar uma solução para $10^6$ 
que será viavel para todos os possiveis valores de $w$. Tendo isso em mente dividimos o a familia em tres secções de $100$ elementos cada,
descritos como:
\begin{align*}
W_1 &= \{ i \cdot 100^0 \mid i \in \mathbb{Z}, \, 1 \le i \le 100 \} \\
W_2 &= \{ i \cdot 100^1 \mid i \in \mathbb{Z}, \, 1 \le i \le 100 \} \\
W_3 &= \{ i \cdot 100^2 \mid i \in \mathbb{Z}, \, 1 \le i \le 100 \}
\end{align*}

Assim para qualquer inteiro $1 \le w \le 10^6$ temos 3 casos:

\begin{itemize}
    \item [a)] se $1 \le w \le 10^2$, para todo $w$ podemos formá-lo utilizando apenas um elemento do conjunto $W_1$.
    
    \item [b)] se $10^2 < w \le 10^4$, podemos formar $w$ somando exatamente dois elementos: um pertencente a $W_1$ 
    representando as duas primeiras casas decimais e outro a representando as ultimas duas $W_2$.
    
    \item [c)] se $10^4 < w \le 10^6$, podemos formar $w$ somando exatamente três elementos distintos: um de $W_1$, 
    um de $W_2$ e um de $W_3$ com o terceiro representando as duas ultimas casas decimais expandindo do ultimo caso.
\end{itemize}

Desse modo, e necessário apenas que para todo $w$ a resposta seja a união entre $W_1$, $W_2$ e $W_3$. 

\section{G - P07 (AtCoder - abc249\_d)}

O problema nos da uma sequência de inteiros $A = (A_1$,\dots,$A_N)$ de tamanho N, e pede como resposta o numero de triplas de indices ($i$,$j$,$k$)
que satisfaçam as seguintes condições:
\begin{itemize}
    \item $1 \le i$,$j$,$k \le N $
    \item $\frac{A_i}{A_j} = A_k$ 
    \item $1 \le N \le 2 \times 10^5$
    \item $1 \le A_i \le 2 \times 10^5 \ (1 \le i \le N)$
\end{itemize} 

\subsection{Solução}

Vamos resolver esse problema usando combinatoria, para isto primeiro precisamos rearranjar todos os $A_i$ da sequencia em um \textit{Heat-map} de tamanho
$K = 2 \times 10^5 + 1$ assim podemos guardar o numero de ocorrencias de cada numero da sequencia e teremos uma variavel \textit{max} que guarda qual é o maior
numero da sequencia, ademais devemos manipular a expressão $\frac{A_i}{A_j} = A_k$ para $A_i = A_j \times A_k$ para facilitar o trabalho de encontrar as triplas. 
Tendo esse \textit{Heat-map} que, relaciona o indice do Array com o numero de ocorrencias daquele índice na sequência, podemos encontrar para todo inteiro $i$ nessa 
sêquencia o numero de triplas de indices que o envolvem e satisfazem as condições estabelecidas, utilizando um \textit{for} sobre todos os indices a partir do i escolhido até que
$i \times j$ seja maior que \textit{max + 1}, multiplicando as ocorrencias dos indices $i$,$j$ e $i \times j$ 
(que nessa contrução são os numeros dos indices correspondentes a $j,k,i$ respectivamente)

\vspace{0.5cm} 

\begin{lstlisting}[language=C++]
for (int i = 1; i < max + 1; i++){
    for (int j = i; j * i < max + 1; j++){
        /* O if abaixo verifica se os indices sao distintos
           e multiplicando a resposta por 2, contando a possivel
           permutacao de indices na multiplicacao da tripla */
        if (i != j) res += (o[i] * o[j] * o[i * j]) * 2; 
        else res += o[i] * o[j] * o[i * j];
    }
}
\end{lstlisting}

\subsection{Análise de Complexidade}

Seja $N = max$. O laço externo itera sobre a variável $i$ de $1$ até $N$. Para cada valor de $i$, o laço interno itera sobre a variável $j$ começando de $i$ enquanto a condição $i \times j \le N$ for verdadeira. Isso implica que, para um dado $i$, o laço interno executa no máximo $\frac{N}{i}$ vezes.

Para determinar o número total de iterações, devemos somar as execuções do laço interno para todos os possíveis valores do laço externo:

$$ \sum_{i=1}^{N} \frac{N}{i} = N \sum_{i=1}^{N} \frac{1}{i} $$

Sabemos que o somatório $\sum_{i=1}^{N} \frac{1}{i}$ representa a série harmônica, cuja soma diverge logaritmicamente. Podemos aproximá-la da seguinte forma:

$$ \sum_{i=1}^{N} \frac{1}{i} \approx \log N $$

Substituindo essa aproximação na nossa equação original, concluímos que o número de operações é proporcional a:

$$ N \log N $$

Portanto, a complexidade de tempo do algoritmo é limitada por $\mathcal{O}(N \log N)$, tornando a execução viável para a densidade de dados apresentados pelo problema.

\end{document}